\section{Design, population, and custody}
\label{sec:design}

\paragraph{Estimand.}
The primary estimand is question-level generalization within a modeled database:
given an established semantic model for a database, how accurately does the
production-governed system answer questions it has not seen during development,
and does the mechanically frozen system generalize to those unseen questions?
The held-out unit is the question, not the
database. All 18 databases appear in both partitions, and full public schema,
column meanings, HKB, and public question text are available during model
construction. This matches the product setting in which a team models its data
once and users then ask an open-ended stream of new questions. It is explicitly
not cold-start performance on an unmodeled database.

A secondary leave-one-database-out study, confined to the development partition,
was specified to estimate whether general HKB-to-semantic-layer transformation
rules transfer to a domain whose questions were excluded from rule development.
It was deferred before execution and is not reported here; the specification is
retained because it fixes the design of any later portability claim. Each fold
would develop rules from 17 databases, apply them to the held-out database's
public schema, column meanings, and HKB without question-driven evidence from
that database, and score that database's development questions, rotating across
all 18, reporting pooled question accuracy and the macro-average across
databases with each fold's question count and modeling effort. If a held-out database receives manual
modeling, the result is transformation portability, not zero-shot database
generalization. This is a development-set study, not a second test set.

\paragraph{Split.}
Eligibility uses only the public \texttt{category} field: 332 \texttt{Query}
records in, 148 \texttt{Management} records out. Algorithm version 1, under
committed seed \texttt{omni-livesqlbench-large-v1-split-v1}, allocates test counts
across \texttt{selected\_database} by Hamilton largest remainder, balances public
\texttt{high\_level} within each database, and orders each stratum by SHA-256 of
seed, database, \texttt{high\_level}, and \texttt{instance\_id}, yielding 231
development and 101 sealed test questions. The 231 development questions are
split again under seed
\texttt{omni-livesqlbench-large-v1-development-split-v1} into 154 \texttt{dev-A}
development questions and 77 \texttt{dev-B} checkpoint questions by the same
database-first, \texttt{high\_level}-balanced procedure. Public
\texttt{conditions} marginals are audited after allocation rather than crossed
into sparse strata; the largest post-allocation deviation is \texttt{dev-B}
\texttt{order}, expected 32 true and 45 false against an actual 37 and 40.

The committed artifacts record source revision, source hash, seed, algorithm
version, record identifiers, and subgroup counts. The eligible manifest hashes to
\hash{ed8a7b3f55e893e05a9e602b34b39509ca29a30798afd63fd9a521f5ffbc7d0e}, the
train identifiers to
\hash{ecdd1e4cd425e41970d3d4e96d7c3cb044cb78d94e34621d320e3d542a5d1c46}, and the
test identifiers to
\hash{c7d4eb1aff3e9be9ffa015adcae96bb9294adf138870db24cdfbed811a9de76f}. Tests
establish that identifiers are unique, train and test are disjoint, their union
is the 332-record eligible manifest, counts are 231 and 101, all 18 databases
appear in both partitions, regeneration from the pinned seeds and source hash is
byte-identical, and no protected field has entered a committed artifact.
Manifests are regenerated, never hand-edited. Membership never changes in
response to any observed outcome.

\paragraph{The executed held-out frame.}
The official Large-v1 Linux loader does not populate two databases. It skips 34
declared tables in \texttt{mental\_healths\_large} and 37 in
\texttt{organ\_transplant\_large}, because the capitalization of those archive
filenames differs from the names the loader requests, so the benchmark's own
reference SQL cannot execute there. Twelve of the 101 sealed questions belong to
those two databases. Before any sealed generation, any label release, and any
outcome access, the held-out frame was amended to the 89 questions on the 16
databases with verified governed deployments, and all four conditions and all
three repetitions used that identical membership, for 1,068 scheduled attempts.
The amendment used no question content, no gold, no annotation, and no
correctness; it used only which databases the public loader can populate. It
narrows the held-out estimand from all 18 benchmark databases to those 16
deployable databases, which is a scope limitation rather than a set of missing
outcomes, and it is reported as such in Section~\ref{sec:threats}. The same
loader defect removes 18 of the 154 development questions from official
scoreability; the development schedule still lists all 154 and reports those 18
as scorer-conformance exclusions, leaving 136 answerable development questions
for the governed arm.

Two further database exclusions apply only to the public direct-SQL development
baseline and not to the held-out frame. \texttt{archeology\_scan\_large}
repeatedly failed to return a usable direct answer across distinct retrieval
settings, and \texttt{cybermarket\_pattern\_large} was excluded after its first
immutable launch failed read-only privilege attestation, even though the external
credential was later repaired. Both were fixed in a committed exclusion artifact
before the affected generations, and their records stay outside the wrong-answer
and missing-row categories.

\paragraph{Information tiers.}
Table~\ref{tab:tiers} states what may be used where. The governing rule is that
hidden training annotations could influence a reusable system under the protocol
but may never become question-specific runtime input. The executed final
candidate does not use that supervision. Development
\texttt{external\_knowledge} identifiers are privileged diagnostic-oracle
metadata: they may be consulted \emph{after} a prediction, to test whether the
required public HKB nodes were represented, whether dependencies compiled, and
whether the system could reach them. They may not select runtime context.
Question-identifier rules, hidden-identifier retrieval hints, and
question-specific gold-derived transformations are prohibited; semantic changes
must be reusable definitions, database-level modeling, transformation rules, or
general harness behavior. If sample queries are used, each must encode a reusable
pattern and not a question-to-gold pair, and material use requires an
ablation.

\begin{table}[htbp]
\centering
\caption{Information tiers.}
\label{tab:tiers}
\begin{tabular}{p{5.1cm}p{2.6cm}p{2.1cm}p{2.4cm}}
\toprule
Information & Development use & Runtime use & Held-out \\
\midrule
Public schema, HKB, column meanings, questions & Allowed & Allowed & Frozen \\
Hidden train SQL, test cases, knowledge identifiers & Offline supervision only & Prohibited & Frozen \\
Hidden test SQL, test cases, knowledge identifiers & Prohibited & Prohibited & Sealed evaluator only \\
\bottomrule
\end{tabular}
\end{table}

\paragraph{Gold custody.}
The complete private attachment stays under human custody outside the repository
and outside any agent-accessible workspace; its SHA-256 is computed without
displaying or parsing its contents. Only after the split commit exists may a
custodian-run extraction emit exactly the 231 training records into an ignored
directory. The release tool verifies the canonical train identifiers and split
metadata against the externally recorded pre-gold commit rather than the mutable
current branch, rejects sources inside the workspace, refuses overwrites, writes
mode \texttt{0600}, and reports only counts and hashes. It never releases the 101
held-out records. Logs may contain identifiers, hashes, counts, status codes, and
aggregate scores, never hidden SQL, hidden knowledge identifiers, or test-case
bodies.

Checkpoint correctness on \texttt{dev-B} sits behind a separate guardian
boundary. The guardian signs exact receipt bytes with a private key held outside
agent scope; the checkpoint capability verifies the detached signature against a
separately recorded public-key digest that is part of the protocol-frozen
configuration and cannot be supplied by a caller, rejects replayed receipt or
output hashes, and irreversibly increments a counter. Each consumption writes
both an immutable checkpoint manifest and a numbered allocation marker, and both
must agree byte-for-byte and be present in the current commit before another
checkpoint is permitted, so marker deletion, manifest deletion, receipt replay,
and local counter reset all fail closed. The protocol maximum is ten
\texttt{dev-B} evaluations, with five to ten intended over the project.

\paragraph{Baseline ordering and the two freezes.}
Before hidden training annotations are released, a public-only baseline is
preserved: transform public schema, column meanings, and HKB into the semantic
layer deterministically; validate the model structurally; freeze the mechanical
transformation revision and model artifacts; generate baseline outputs on all 231
development questions; and only then release train annotations for scoring and
failure analysis. This separates initial transformation quality from any
improvement that might have been obtained through supervision. The baseline may be scored once its immutable
outputs exist, but may never be regenerated retrospectively. The protocol
permitted supervision on development questions, but the executed final system is
the frozen mechanical baseline and receives no question-level supervision. It is
still not a cold-start database evaluation because public schema, column
meanings, and HKB are compiled into the model.

\emph{Freeze A}, the pre-gold protocol freeze, commits the eligible population,
both splits, custody and information rules, the C1 through C4 definitions, both
scoring policies, repetition, order, and failure rules, endpoints and statistics,
unblinding and adjudication policy, the ledger schema, and the guardian
public-key pin. Its commit hash is recorded before any hidden development label
is released, and the control plane verifies configuration and manifests against
that recorded commit, not the working tree. Freeze A deliberately does not
freeze the final system. \emph{Freeze B}, the pre-test system freeze, occurs
after the recorded decision to cut optimization and consume no checkpoints, and before any
test generation: it commits and hashes the semantic models and transformation
code, all four harness configurations and prompts, retrieval configuration and
allowed tools, model and provider information and budgets, retry and
failure-classification policy, database snapshot identifiers, the trial schedule,
and the evaluator version. All 1,068 scheduled generations then complete, and
only then are they scored. No held-out result may feed a model, harness, prompt,
retrieval, or retry change.

Both freezes executed. The untuned system that produced the held-out generations
is frozen at successor commit
\hash{8b0c7393d564d9ecc2c2f84ba7446d610c1a0a6d} with direct-child control commit
\hash{94cc0d9483c944d7dc13ed651c8fc2ef077f33ab} and Freeze-B manifest SHA-256
\hash{e1c9f1967422822c848c18a17ba759d4e4fbc7f21aa0fe3ae1045b9236ae4730}. Scoring
ran from correction-forward system
\hash{0a5aee423b4a0d5bb396b3d9764f8e9e24f31254} against control
\hash{fe4660df8dacdca07da310ddfda4158b82895ba9} under Freeze-B v10 SHA-256
\hash{ff10083bf70d82bd483d12e98751d9bf7f5d4236c42fac3ba921405d87953a05},
preserving the generation binding above. The correctness-free scoring receipt
hashes to
\hash{534e28b954d4d13dfdd9100fc6a184fba3eb3720d8cd7cf7d43c92713cb258f7}. No
optimized held-out candidate was constructed after the sealed result.
